\section{Part H. 왜 복잡도인가?}

% 31
\begin{frame}{복잡도는 계산 비용의 함수다}
  \begin{columns}[T,onlytextwidth]
    \column{.47\textwidth}
      \begin{block}{Time complexity}
        입력 크기 $n$에 따라 필요한 기본 연산 수는 어떻게 변하는가?
      \end{block}
    \column{.47\textwidth}
      \begin{block}{Space complexity}
        입력 외에 필요한 추가 메모리는 어떻게 변하는가?
      \end{block}
  \end{columns}
  \vfill
  \begin{alertblock}{오늘의 비용 모델}
    비교, 대입, 산술 같은 기본 연산 하나는 상수 시간이라고 가정한다.
  \end{alertblock}
\end{frame}

% 32
\begin{frame}{실행 시간만 재면 충분할까?}
  두 프로그램의 예측 시간이 다음과 같다고 하자.
  \[
    T_A(n)=30n+8,\qquad T_B(n)=n^2+1 \quad (\mu s)
  \]
  \begin{columns}[T,onlytextwidth]
    \column{.47\textwidth}
      \textbf{Benchmark}\par
      특정 기계·언어·입력에서 실제 성능을 보여 준다.
    \column{.47\textwidth}
      \textbf{Analysis}\par
      구현 세부를 넘어 입력 규모에 따른 추세를 설명한다.
  \end{columns}
  \vfill
  수백만 개를 처리한다면 $T_A$를 선택한다. 작은 $n$의 승자가 큰 $n$에서도 이긴다는 보장은 없다.
\end{frame}

% 33
\begin{frame}{비용은 입력에 따라 달라진다}
  같은 크기 $n$에서도 데이터 배치가 다르면 수행 연산 수가 달라진다.
  \begin{description}
    \item[최선 경우(Best case)] 가장 유리한 입력의 비용
    \item[평균 경우(Average case)] 입력 분포를 가정한 기대 비용
    \item[최악 경우(Worst case)] 가장 불리한 입력의 비용
  \end{description}
  \begin{block}{왜 worst case를 자주 쓰나?}
    모든 크기-$n$ 입력에 대한 상한을 주며, 별도 확률 분포가 필요 없다.
  \end{block}
\end{frame}

% 34
\begin{frame}{복잡도 정리}
  분석을 시작할 때 네 가지를 먼저 적는다.
  \begin{enumerate}
    \item 입력 크기 $n$은 무엇인가?
    \item 어떤 연산을 셀 것인가?
    \item best, average, worst 중 무엇인가?
    \item 시간과 공간 중 무엇을 분석하는가?
  \end{enumerate}
  \begin{takeaway}
    복잡도는 초 단위 측정값이 아니라, \textbf{입력 크기에 따른 자원 증가 함수}다.
  \end{takeaway}
\end{frame}
